\documentclass[11pt,a4paper]{article}
\usepackage{kotex} % For Korean
\usepackage[margin=25mm]{geometry} % Page margins
\usepackage{hyperref} % For clickable links and PDF metadata
\usepackage{xcolor} % For colors
\usepackage[most]{tcolorbox} % For colored boxes
\usepackage{booktabs} % For professional tables
\usepackage{graphicx} % For images (though none in transcript, good to have)
\usepackage{amsmath, amssymb} % For math symbols
\usepackage{listings} % For code blocks
\usepackage{setspace} % For line spacing
\usepackage{fancyhdr} % For headers and footers
\usepackage{adjustbox} % For resizing content, especially tables

% Line spacing
\onehalfspacing
% Paragraph settings
\setlength{\parskip}{0.6em}
\setlength{\parindent}{0pt}
% Table row stretch
\renewcommand{\arraystretch}{1.2}

% Listing settings for code blocks
\lstset{
  basicstyle=\ttfamily\small,
  breaklines=true,
  numbers=left,
  numbersep=6pt,
  frame=single,
  captionpos=b, % Caption at bottom
  keywordstyle=\color{blue},
  stringstyle=\color{red},
  commentstyle=\color{green!50!black},
  showstringspaces=false,
  tabsize=2,
  extendedchars=true,
  inputencoding=utf8/latin1,
  literate={á}{{\'a}}1 {é}{{\'e}}1 {í}{{\'i}}1 {ó}{{\'o}}1 {ú}{{\'u}}1
           {Á}{{\'A}}1 {É}{{\'E}}1 {Í}{{\'I}}1 {Ó}{{\'O}}1 {Ú}{{\'U}}1
           {à}{{\`a}}1 {è}{{\`e}}1 {ì}{{\`i}}1 {ò}{{\`o}}1 {ù}{{\`u}}1
           {À}{{\`A}}1 {È}{{\`E}}1 {Ì}{{\`I}}1 {Ò}{{\`O}}1 {Ù}{{\`U}}1
           {ä}{{\"a}}1 {ë}{{\"e}}1 {ï}{{\"i}}1 {ö}{{\"o}}1 {ü}{{\"u}}1
           {Ä}{{\"A}}1 {Ë}{{\"E}}1 {Ï}{{\"I}}1 {Ö}{{\"O}}1 {Ü}{{\"U}}1
           {ç}{{\c c}}1 {Ç}{{\c C}}1 {ñ}{{\~n}}1 {Ñ}{{\~N}}1
           {ß}{{\ss}}1 {œ}{{\oe}}1 {Œ}{{\OE}}1 {æ}{{\ae}}1 {Æ}{{\AE}}1
           {ø}{{\o}}1 {Ø}{{\O}}1 {å}{{\aa}}1 {Å}{{\AA}}1
           {€}{{\euro}}1 {£}{{\pounds}}1 {¥}{{\yen}}1 {₩}{{\won}}1
           {ㄱ}{{\textHangul{ㄱ}}}1 {ㄴ}{{\textHangul{ㄴ}}}1 {ㄷ}{{\textHangul{ㄷ}}}1 {ㄹ}{{\textHangul{ㄹ}}}1 {ㅁ}{{\textHangul{ㅁ}}}1 {ㅂ}{{\textHangul{ㅂ}}}1 {ㅅ}{{\textHangul{ㅅ}}}1 {ㅇ}{{\textHangul{ㅇ}}}1 {ㅈ}{{\textHangul{ㅈ}}}1 {ㅊ}{{\textHangul{ㅊ}}}1 {ㅋ}{{\textHangul{ㅋ}}}1 {ㅌ}{{\textHangul{ㅌ}}}1 {ㅍ}{{\textHangul{ㅍ}}}1 {ㅎ}{{\textHangul{ㅎ}}}1
           {ㅏ}{{\textHangul{ㅏ}}}1 {ㅑ}{{\textHangul{ㅑ}}}1 {ㅓ}{{\textHangul{ㅓ}}}1 {ㅕ}{{\textHangul{ㅕ}}}1 {ㅗ}{{\textHangul{ㅗ}}}1 {ㅛ}{{\textHangul{ㅛ}}}1 {ㅜ}{{\textHangul{ㅜ}}}1 {ㅠ}{{\textHangul{ㅠ}}}1 {ㅡ}{{\textHangul{ㅡ}}}1 {ㅣ}{{\textHangul{ㅣ}}}1
           {가}{{\textHangul{가}}}1 {나}{{\textHangul{나}}}1 {다}{{\textHangul{다}}}1 {라}{{\textHangul{라}}}1 {마}{{\textHangul{마}}}1 {바}{{\textHangul{바}}}1 {사}{{\textHangul{사}}}1 {아}{{\textHangul{아}}}1 {자}{{\textHangul{자}}}1 {차}{{\textHangul{차}}}1 {카}{{\textHangul{카}}}1 {타}{{\textHangul{타}}}1 {파}{{\textHangul{파}}}1 {하}{{\textHangul{하}}}1
}

% Hyperref settings
\hypersetup{
  colorlinks=true,
  linkcolor=blue,
  citecolor=green,
  urlcolor=cyan,
  pdftitle={CSCI103 Lecture 1: 데이터 엔지니어링 소개},
  pdfauthor={UIUC Student},
  pdfsubject={CSCI E-103: Introduction to the Intellectual Enterprises of CS}
}

% Tcolorbox definitions for various box types
\definecolor{myblue}{RGB}{220,230,241}
\definecolor{mygreen}{RGB}{220,241,230}
\definecolor{myred}{RGB}{241,220,220}
\definecolor{myyellow}{RGB}{255,255,204}

\newtcolorbox{summarybox}[1][]{
  colback=myblue,
  colframe=blue!75!black,
  fonttitle=\bfseries,
  title={요약},
  #1
}

\newtcolorbox{notebox}[1][]{
  colback=myyellow,
  colframe=orange!75!black,
  fonttitle=\bfseries,
  title={참고},
  #1
}

\newtcolorbox{warningbox}[1][]{
  colback=myred,
  colframe=red!75!black,
  fonttitle=\bfseries,
  title={주의},
  #1
}

\newtcolorbox{examplebox}[1][]{
  colback=mygreen,
  colframe=green!75!black,
  fonttitle=\bfseries,
  title={예시},
  #1
}

% Header and Footer
\pagestyle{fancy}
\fancyhf{} % Clear all header and footer fields
\fancyhead[L]{CSCI103 Lecture 1: 데이터 엔지니어링 소개}
\fancyhead[R]{\today}
\fancyfoot[C]{\thepage}
\renewcommand{\headrulewidth}{0.4pt}
\renewcommand{\footrulewidth}{0.4pt}

\begin{document}

% Title page (optional, but good for a formal document)
\begin{titlepage}
    \centering
    \vspace*{1cm}
    {\Huge\bfseries CSCI103 Lecture 1: 데이터 엔지니어링 소개\par}
    \vspace{1.5cm}
    {\Large\bfseries CSCI E-103: Introduction to the Intellectual Enterprises of CS\par}
    \vspace{2cm}
    {\Large\itshape 수업자료 통합 노트\par}
    \vfill
    {\large UIUC Student\par}
    \vspace{1cm}
    {\large \today\par}
\end{titlepage}

\newpage

% --- 개요 (Overview) ---
\section*{개요}
\addcontentsline{toc}{section}{개요}
본 문서는 \texttt{CSCI103 Lecture 1} 강의 자료를 통합하여 작성된 시험 대비용 LaTeX 노트입니다.
데이터 엔지니어링의 기본 개념, 빅데이터의 특성, 데이터 플랫폼의 진화 과정, 그리고 클라우드 기반 실습 환경인 Databricks 사용법을 다룹니다.
강의 목표는 비즈니스 문제 해결을 위한 데이터 엔지니어링의 중요성을 이해하고, 실제 분석 환경에서 데이터를 효과적으로 다루는 기술을 습득하는 것입니다.
주요 평가 포인트는 과제, 퀴즈, 최종 그룹 프로젝트 발표이며, 적극적인 참여가 중요합니다.
이 문서는 초심자도 완벽히 이해할 수 있도록 상세한 설명과 예시를 포함합니다.

\newpage

% --- 용어 정리 (Glossary) ---
\section*{용어 정리}
\addcontentsline{toc}{section}{용어 정리}
다음은 강의에서 소개된 주요 용어들을 쉽게 설명한 표입니다.

\begin{adjustbox}{max width=\textwidth}
\begin{tabular}{llll}
\toprule
\textbf{용어} & \textbf{쉬운 설명} & \textbf{원어} & \textbf{비고} \\
\midrule
데이터 엔지니어링 & 원시 데이터를 가치 있는 통찰로 바꾸는 과정 & Data Engineering & 데이터 파이프라인 구축 및 관리 \\
빅데이터 & 기존 도구로 처리하기 어려운 대규모 데이터 & Big Data & 5V(Volume, Velocity, Variety, Veracity, Value) 특성 \\
Databricks & 빅데이터 처리 및 머신러닝을 위한 통합 플랫폼 & Databricks & Spark 기반, 레이크하우스 아키텍처 \\
데이터 페르소나 & 데이터 관련 업무를 수행하는 다양한 역할 & Data Personas & 엔지니어, 과학자, 분석가 등 \\
분산 컴퓨팅 & 여러 컴퓨터가 협력하여 작업을 처리하는 방식 & Distributed Computing & 빅데이터 처리에 필수적 \\
머신러닝 (ML) & 데이터 학습을 통해 예측/결정하는 인공지능 분야 & Machine Learning & 데이터 엔지니어링의 주요 결과물 \\
ETL & 데이터 추출, 변환, 적재 과정 & Extract, Transform, Load & 데이터 파이프라인의 핵심 \\
데이터 파이프라인 & 데이터가 흐르고 처리되는 자동화된 경로 & Data Pipeline & 원시 데이터 $\rightarrow$ 정제 데이터 \\
레이크하우스 & 데이터 레이크와 데이터 웨어하우스의 장점 결합 & Lakehouse & Databricks의 핵심 아키텍처 \\
ACID & 데이터베이스 트랜잭션의 4가지 속성 & Atomicity, Consistency, Isolation, Durability & 관계형 데이터베이스의 신뢰성 보장 \\
BASE & NoSQL 데이터베이스의 3가지 속성 & Basically Available, Soft state, Eventual consistency & 유연성 및 확장성 중시 \\
CAP 이론 & 분산 시스템의 3가지 속성 중 2가지만 만족 가능 & Consistency, Availability, Partition tolerance & 시스템 설계 시 트레이드오프 \\
OLTP & 실시간 트랜잭션 처리 시스템 & Online Transactional Processing & 은행 거래, 온라인 쇼핑 등 \\
OLAP & 대규모 데이터 분석 시스템 & Online Analytical Processing & 보고서, 비즈니스 인텔리전스 \\
Hadoop & 대규모 데이터 분산 처리를 위한 오픈소스 프레임워크 & Hadoop & 데이터 레이크 시대의 시작 \\
Spark & Hadoop의 한계를 극복한 인메모리 분산 처리 엔진 & Spark & 빠른 속도, 다양한 기능 지원 \\
RDD & Spark의 기본 데이터 추상화 단위 & Resilient Distributed Dataset & 불변성, 분산 처리 \\
데이터프레임 & Spark의 구조화된 데이터 처리 API & DataFrame & RDD보다 사용하기 쉬운 추상화 \\
IaaS & 인프라를 서비스로 제공하는 클라우드 모델 & Infrastructure as a Service & 가상 머신, 스토리지 등 \\
PaaS & 플랫폼을 서비스로 제공하는 클라우드 모델 & Platform as a Service & 개발 환경, 데이터베이스 등 \\
SaaS & 소프트웨어를 서비스로 제공하는 클라우드 모델 & Software as a Service & Gmail, Salesforce, Databricks 등 \\
Unity Catalog & Databricks의 통합 데이터 거버넌스 솔루션 & Unity Catalog & 데이터 접근 제어, 메타데이터 관리 \\
\bottomrule
\end{tabular}
\end{adjustbox}

\newpage

% --- 핵심 개념·원리 (Core Concepts & Principles) ---
\section*{핵심 개념·원리}
\addcontentsline{toc}{section}{핵심 개념·원리}

\subsection*{1. 데이터 엔지니어링의 정의와 중요성}

\begin{summarybox}
데이터 엔지니어링은 원시 데이터를 수집, 변환, 관리하여 비즈니스 통찰력을 얻고 머신러닝 모델을 구축하는 데 사용할 수 있도록 준비하는 과정입니다. 이는 현대 비즈니스에서 데이터의 가치를 극대화하는 핵심적인 역할입니다.
\end{summarybox}

\textbf{왜 필요한가?}
오늘날 대부분의 비즈니스는 데이터를 기반으로 운영됩니다. 하지만 데이터는 종종 여러 소스에 흩어져 있고, 정제되지 않은 '날것'의 형태입니다. 이러한 원시 데이터는 그 자체로는 가치가 낮습니다. 데이터 엔지니어링은 이 원시 데이터를 분석 가능한 형태로 가공하고, 신뢰할 수 있는 파이프라인을 통해 지속적으로 제공함으로써 비즈니스 의사결정과 혁신을 가능하게 합니다.

\textbf{직관적 비유:}
데이터 엔지니어는 마치 정수기 필터와 배관공을 합쳐놓은 것과 같습니다.
수도꼭지에서 나오는 흙탕물(원시 데이터)을 바로 마실 수는 없습니다.
데이터 엔지니어는 이 흙탕물을 깨끗하게 정수(변환)하고, 각 가정(분석가, 데이터 과학자)에 필요한 형태로 안정적으로 공급(파이프라인)하는 역할을 합니다. 그래야 사람들이 깨끗한 물(통찰력)을 마시고 건강(비즈니스 성장)해질 수 있습니다.

\textbf{기술적 설명:}
데이터 엔지니어링은 소프트웨어 엔지니어링, 데이터 과학, 비즈니스 인텔리전스(BI)의 교차점에 있습니다.
주요 업무는 다음과 같습니다:
\begin{itemize}
    \item \textbf{데이터 구조 및 분산 컴퓨팅:} 대규모 데이터를 효율적으로 저장하고 처리하기 위한 구조를 설계하고 분산 시스템을 활용합니다.
    \item \textbf{프로그래밍 및 도구:} Python, Scala, SQL 등 다양한 언어와 Spark, Databricks 같은 도구를 사용하여 데이터를 처리합니다.
    \item \textbf{ETL 및 데이터 파이프라인 구축:} 데이터를 추출(Extract), 변환(Transform), 적재(Load)하는 과정을 자동화하는 파이프라인을 설계하고 구현합니다.
    \item \textbf{소프트웨어 엔지니어링 원칙 적용:} 설계, 개발, 웹/모바일 앱, DevOps, 서비스 배포 등 소프트웨어 개발의 모범 사례를 따릅니다.
    \item \textbf{데이터 과학 및 비즈니스 인텔리전스 이해:} 머신러닝 모델링, 알고리즘, 비즈니스 요구사항을 이해하여 데이터가 최종적으로 어떻게 활용될지 고려합니다.
\end{itemize}

\subsection*{2. 데이터 관련 페르소나 (Data Personas)}

데이터를 다루는 다양한 역할들이 있으며, 이들은 서로 협력하여 데이터의 가치를 창출합니다. 데이터 엔지니어는 이 모든 역할과 긴밀하게 연결되어 있습니다.

\begin{adjustbox}{max width=\textwidth}
\begin{tabular}{lll}
\toprule
\textbf{페르소나} & \textbf{주요 역할} & \textbf{데이터 엔지니어와의 관계} \\
\midrule
\textbf{데이터 엔지니어} & 데이터 파이프라인 구축 및 유지보수, 데이터 변환 및 준비 & 모든 역할의 기반 데이터 제공 \\
\textbf{BI 분석가} & 정제된 데이터를 분석하여 대시보드 및 보고서 생성 & 엔지니어가 준비한 데이터를 소비 \\
\textbf{데이터 과학자/ML 실무자} & 머신러닝 모델 개발, 피처 엔지니어링, 데이터 기반 예측 & 엔지니어가 준비한 고품질 데이터를 활용하여 모델 학습 \\
\textbf{DevOps/MLOps} & 데이터 파이프라인 및 ML 모델 배포 및 자동화 지원 & 엔지니어가 구축한 파이프라인의 운영 효율성 담당 \\
\textbf{데이터 리더} & 데이터 전략 수립, IT 조직의 데이터 방향성 제시 & 엔지니어에게 비즈니스 요구사항 전달, 데이터 활용 가이드 \\
\bottomrule
\end{tabular}
\end{adjustbox}

\begin{notebox}
최종 프로젝트는 그룹 프로젝트로 진행되며, 각 팀원은 이러한 데이터 페르소나 중 하나를 맡아 프로젝트의 특정 측면을 발표하게 됩니다. 이는 실제 업무 환경을 반영한 것입니다.
\end{notebox}

\subsection*{3. 빅데이터의 특성: 5V}

오늘날 대부분의 데이터셋은 '빅데이터'로 분류되며, 그 특성은 5가지 'V'로 설명됩니다.

\begin{adjustbox}{max width=\textwidth}
\begin{tabular}{lll}
\toprule
\textbf{특성} & \textbf{설명} & \textbf{예시} \\
\midrule
\textbf{Volume (규모)} & 데이터의 물리적 크기 (테라바이트, 페타바이트 이상) & 수십억 건의 고객 거래 기록, 수천 시간의 비디오 데이터 \\
\textbf{Velocity (속도)} & 데이터가 생성되고 처리되는 속도 (실시간, 배치) & 주식 시장 데이터 (초당 수천 건), IoT 센서 데이터 (지속적인 스트림) \\
\textbf{Variety (다양성)} & 데이터의 형태 (정형, 반정형, 비정형) & 관계형 데이터베이스(정형), JSON/XML(반정형), 오디오/비디오/이미지(비정형) \\
\textbf{Veracity (정확성)} & 데이터의 신뢰성 및 정확성 & 잘못된 센서 값, 중복된 고객 정보, 편향된 설문 응답 \\
\textbf{Value (가치)} & 데이터를 통해 얻을 수 있는 잠재적 비즈니스 가치 & 고객 행동 분석을 통한 맞춤형 추천, 질병 예측 모델 \\
\bottomrule
\end{tabular}
\end{adjustbox}

\begin{examplebox}
\textbf{5V 예시: 소셜 미디어 데이터}
\begin{itemize}
    \item \textbf{Volume:} 매일 수십 페타바이트의 게시물, 좋아요, 댓글이 생성됩니다.
    \item \textbf{Velocity:} 사용자들이 실시간으로 콘텐츠를 생성하고 상호작용합니다.
    \item \textbf{Variety:} 텍스트(게시물), 이미지, 비디오(비정형), 사용자 프로필(정형), API 응답(반정형) 등 다양한 형태입니다.
    \item \textbf{Veracity:} 가짜 뉴스, 봇 활동, 스팸 등으로 인해 데이터의 신뢰성이 떨어질 수 있습니다.
    \item \textbf{Value:} 사용자 관심사 분석, 트렌드 파악, 광고 타겟팅, 여론 분석 등 막대한 비즈니스 가치를 가집니다.
\end{itemize}
\end{examplebox}

\subsection*{4. 데이터 분류 기준}

데이터는 다양한 기준에 따라 분류될 수 있으며, 이는 데이터 처리 전략을 결정하는 데 중요합니다.

\begin{adjustbox}{max width=\textwidth}
\begin{tabular}{lll}
\toprule
\textbf{분류 기준} & \textbf{설명} & \textbf{예시} \\
\midrule
\textbf{크기} &
\begin{itemize}
    \item \textbf{빅데이터:} 테라바이트 이상, 분산 컴퓨팅 필요
    \item \textbf{일반 데이터:} 단일 컴퓨터로 처리 가능
\end{itemize} &
\begin{itemize}
    \item 수십억 건의 웹 로그 (빅데이터)
    \item 개인 노트북에 저장된 스프레드시트 (일반 데이터)
\end{itemize} \\
\textbf{구조} &
\begin{itemize}
    \item \textbf{정형 데이터 (Structured):} 스키마가 명확
    \item \textbf{반정형 데이터 (Semi-structured):} 스키마가 데이터 내에 포함
    \item \textbf{비정형 데이터 (Unstructured):} 스키마 없음
\end{itemize} &
\begin{itemize}
    \item 관계형 데이터베이스 (정형)
    \item JSON, XML (반정형)
    \item 이미지, 오디오, 비디오, 문서 (비정형)
\end{itemize} \\
\textbf{속도} &
\begin{itemize}
    \item \textbf{배치 (Batch):} 일정 시간 동안 모아서 일괄 처리
    \item \textbf{스트리밍 (Streaming):} 데이터가 생성되는 즉시 연속적으로 처리
    \item \textbf{마이크로 배치 (Micro-batch):} 스트리밍과 배치 중간 형태
\end{itemize} &
\begin{itemize}
    \item 일일 재고 보고서 (배치)
    \item 주식 시세, IoT 센서 데이터 (스트리밍)
\end{itemize} \\
\textbf{변경 빈도} &
\begin{itemize}
    \item \textbf{거의 변경 없음:} 매우 드물게 업데이트
    \item \textbf{가끔 변경:} 주기적으로 업데이트
    \item \textbf{자주 변경:} 빈번하게 업데이트
    \item \textbf{매우 자주 변경:} 지속적으로 업데이트
\end{itemize} &
\begin{itemize}
    \item 인구 조사 데이터 (거의 변경 없음)
    \item 공장 기계 상태 (가끔 변경)
    \item 소셜 미디어 상호작용 (자주 변경)
    \item 날씨 데이터 (매우 자주 변경)
\end{itemize} \\
\bottomrule
\end{tabular}
\end{adjustbox}

\subsection*{5. 데이터 온도와 처리 속도}

데이터는 그 중요성과 긴급성에 따라 '온도' 개념으로 분류되며, 이에 따라 적절한 처리 속도가 요구됩니다.

\begin{adjustbox}{max width=\textwidth}
\begin{tabular}{lll}
\toprule
\textbf{데이터 온도} & \textbf{설명} & \textbf{처리 속도} \\
\midrule
\textbf{Hot Data (뜨거운 데이터)} & 최근에 생성되었고, 빠르게 변화하며, 즉각적인 분석이 필요한 데이터 & \textbf{실시간 (Real-time):} 데이터 도착 즉시 처리 \\
\textbf{Warm Data (따뜻한 데이터)} & 비교적 최근 데이터지만, 즉각적인 처리가 필수는 아닌 데이터 & \textbf{인터랙티브 (Interactive):} 필요 시 즉시 쿼리 및 분석 \\
\textbf{Cold Data (차가운 데이터)} & 오래되었고, 거의 변경되지 않으며, 보관 및 가끔의 분석에 사용되는 데이터 & \textbf{배치 (Batch):} 일정 시간 모아서 일괄 처리 \\
\bottomrule
\end{tabular}
\end{adjustbox}

\begin{examplebox}
\textbf{데이터 온도 및 처리 속도 예시:}
\begin{itemize}
    \item \textbf{Hot Data (실시간):} 온라인 사기 탐지 시스템. 거래가 발생하는 즉시 분석하여 사기 여부를 판단해야 합니다. 5분, 15분 지연은 큰 손실로 이어질 수 있습니다.
    \item \textbf{Warm Data (인터랙티브):} 고객 서비스 상담원이 고객의 최근 1주일 구매 내역을 조회하여 상담에 활용하는 경우. 즉각적인 응답이 필요하지만, 초 단위의 실시간은 아닙니다.
    \item \textbf{Cold Data (배치):} 연간 재무 보고서 작성. 지난 1년간의 모든 거래 데이터를 밤새 일괄 처리하여 다음 날 아침에 보고서를 생성합니다.
\end{itemize}
\end{examplebox}

\begin{warningbox}
비즈니스에서 항상 '실시간' 처리를 요구하지만, 모든 데이터가 실시간으로 처리될 필요는 없습니다. 실시간 처리는 비용이 많이 들기 때문에, 해당 통찰이 정말로 '실행 가능(actionable)'하고 비즈니스 가치가 높은지 신중하게 판단해야 합니다.
\end{warningbox}

\subsection*{6. 데이터 가치 증대 과정 (Refinement)}

원시 데이터는 여러 단계를 거쳐 정제되고 변환될수록 비즈니스 가치가 높아집니다.

\begin{figure}[h!]
    \centering
    % Placeholder for the image described in the transcript
    % The transcript describes a visual LinkedIn post showing raw data as a messy pile,
    % then sorted/arranged data, then aggregated data, and finally a diamond.
    % Since I cannot generate images, I will describe it.
    \begin{tcolorbox}[colback=gray!10, colframe=gray!50!black, title={데이터 정제 과정 (개념도)}]
        \centering
        \textbf{원시 데이터 (Raw Data)} \\
        (무질서한 더미) \\
        $\downarrow$ \\
        \textbf{정렬/정돈된 데이터 (Sorted/Arranged Data)} \\
        (일부 질서가 잡힘) \\
        $\downarrow$ \\
        \textbf{집계된 데이터 (Aggregated Data)} \\
        (패턴과 요약이 보임) \\
        $\downarrow$ \\
        \textbf{통찰력/가치 (Insights/Value)} \\
        (다이아몬드처럼 빛나는 최종 결과물)
    \end{tcolorbox}
    \caption{원시 데이터에서 가치 있는 통찰로의 변환 과정}
    \label{fig:data_refinement}
\end{figure}

\begin{itemize}
    \item \textbf{원시 데이터 (Raw Data):} 소스 시스템에서 직접 추출된, 가공되지 않은 데이터. (예: 흙탕물)
    \item \textbf{정제 및 변환 (Cleansing & Transformation):} 데이터의 오류를 수정하고, 형식을 표준화하며, 필요한 형태로 가공합니다. (예: 정수 필터링)
    \item \textbf{집계 (Aggregation):} 데이터를 요약하고 그룹화하여 패턴을 찾습니다. (예: 물을 병에 담아 분류)
    \item \textbf{통찰력 (Insights):} 정제된 데이터를 분석하여 비즈니스 의사결정에 도움이 되는 의미 있는 정보를 도출합니다. (예: 깨끗한 생수)
\end{itemize}

\begin{notebox}
Databricks에서는 이러한 데이터 정제 과정을 \textbf{메달리온 아키텍처(Medallion Architecture)}로 설명합니다.
\begin{itemize}
    \item \textbf{Bronze (브론즈):} 원시 데이터 (최소한의 정제)
    \item \textbf{Silver (실버):} 정제되고 필터링된 데이터 (비즈니스 규칙 적용)
    \item \textbf{Gold (골드):} 집계되고 분석 준비가 완료된 데이터 (최종 사용자용)
\end{itemize}
이 아키텍처는 데이터의 품질과 신뢰성을 단계적으로 높여줍니다.
\end{notebox}

\subsection*{7. 데이터 플랫폼의 진화}

데이터를 다루는 기술 스택은 시대에 따라 계속 발전해 왔지만, 핵심적인 사용 사례(use cases)와 데이터 자체의 가치는 변하지 않습니다.

\begin{adjustbox}{max width=\textwidth}
\begin{tabular}{llll}
\toprule
\textbf{시대} & \textbf{주요 기술/개념} & \textbf{특징} & \textbf{한계/변화 요인} \\
\midrule
\textbf{1960년대} & DBMS (데이터베이스 관리 시스템) & 플랫 파일에서 관계형 데이터베이스로 전환 & 데이터 볼륨 증가, 비정형 데이터 미지원 \\
\textbf{1980년대} & 데이터 웨어하우스 (Data Warehouse) & 차원 모델링, 데이터 마트, MPP (대규모 병렬 처리) & 고비용, 독점 시스템, BI 중심, 비정형 데이터 미지원 \\
\textbf{2000년대} & 웹 및 비정형 데이터 & 웹의 부상, 오디오/비디오/이미지, 메타데이터 중요성 증가 & SQL의 한계, NoSQL 등장, 스트리밍 초기 단계 \\
\textbf{2010년대 초} & Hadoop (하둡) & 오픈소스, 저비용 분산 처리, 데이터 레이크 개념 도입 & 데이터 품질/신뢰성 문제, 복잡한 관리 \\
\textbf{2010년대 중반} & Spark (스파크) & Hadoop의 한계 극복, 인메모리 처리, 빠른 속도, 클라우드 채택 & 클라우드 컴퓨팅 (AWS, Azure, GCP) 확산, 스토리지/컴퓨트 분리 \\
\textbf{2020년대} & 레이크하우스 (Lakehouse) & 데이터 레이크와 데이터 웨어하우스의 장점 결합 & 데이터 메시, 데이터 패브릭 등 새로운 아키텍처 등장 \\
\bottomrule
\end{tabular}
\end{adjustbox}

\begin{notebox}
\textbf{기술 부채 (Tech Debt):} 오래된 기술 스택을 계속 사용하면 유지보수 비용이 증가하고 새로운 기능을 도입하기 어려워지는 현상입니다. 비즈니스는 결국 '총알을 물고(bite the bullet)' 더 현대적인 스택으로 전환해야 합니다.
\end{notebox}

\subsection*{8. SQL 기반 시스템 vs NoSQL 기반 시스템}

데이터베이스는 크게 SQL(관계형)과 NoSQL(비관계형)로 나눌 수 있으며, 각각 다른 특성과 사용 사례를 가집니다.

\subsubsection*{8.1. SQL 기반 시스템 (ACID)}

\begin{summarybox}
SQL 기반 시스템은 관계형 데이터베이스의 핵심으로, 데이터의 일관성과 신뢰성을 최우선으로 합니다. \textbf{ACID} 속성을 통해 트랜잭션의 무결성을 보장합니다.
\end{summarybox}

\textbf{ACID 속성:}
\begin{itemize}
    \item \textbf{Atomicity (원자성):} 트랜잭션 내의 모든 작업은 완전히 성공하거나, 아니면 완전히 실패하여 이전 상태로 되돌려집니다 (All or Nothing).
    \begin{examplebox}
    \textbf{예시:} 은행 계좌 이체. A 계좌에서 돈을 인출하고 B 계좌로 입금하는 두 작업이 모두 성공해야 합니다. 만약 인출만 되고 입금이 실패하면, 전체 트랜잭션은 롤백되어 A 계좌의 돈도 다시 돌아옵니다.
    \end{examplebox}
    \item \textbf{Consistency (일관성):} 트랜잭션이 성공적으로 완료되면 데이터베이스는 항상 유효한 상태를 유지합니다. 데이터베이스 규칙(제약 조건)을 위반하지 않습니다.
    \begin{examplebox}
    \textbf{예시:} 재고 관리 시스템. 상품이 판매되면 재고 수량이 정확히 감소해야 합니다. 10개 팔렸는데 9개만 감소하면 일관성이 깨집니다.
    \end{examplebox}
    \item \textbf{Isolation (격리성):} 여러 트랜잭션이 동시에 실행될 때, 각 트랜잭션은 다른 트랜잭션의 영향을 받지 않고 독립적으로 실행되는 것처럼 보입니다.
    \begin{examplebox}
    \textbf{예시:} 두 고객이 동시에 같은 상품을 구매하려고 할 때, 한 고객의 구매가 완료될 때까지 다른 고객은 해당 상품의 재고 상태를 변경할 수 없습니다.
    \end{examplebox}
    \item \textbf{Durability (영속성):} 트랜잭션이 성공적으로 완료되면, 그 결과는 시스템 오류(정전 등)에도 불구하고 영구적으로 저장됩니다.
    \begin{examplebox}
    \textbf{예시:} 결제 완료 후 시스템이 다운되더라도, 결제 기록은 사라지지 않고 유지됩니다.
    \end{examplebox}
\end{itemize}

\textbf{주요 사용 사례:}
\begin{itemize}
    \item 고도로 구조화된 데이터 (정형 데이터)
    \item 예측 가능한 입력 및 출력
    \item 금융 시스템, 회계 장부, 재고 관리 등 데이터의 정확성과 일관성이 매우 중요한 경우
\end{itemize}

\subsubsection*{8.2. NoSQL 기반 시스템 (BASE)}

\begin{summarybox}
NoSQL 기반 시스템은 대규모 분산 환경에서 유연성, 확장성, 고가용성을 중시합니다. \textbf{BASE} 속성을 통해 일관성보다는 가용성과 성능에 초점을 맞춥니다.
\end{summarybox}

\textbf{BASE 속성:}
\begin{itemize}
    \item \textbf{Basically Available (기본적 가용성):} 시스템의 일부 노드에 장애가 발생하더라도 전체 시스템은 계속해서 작동하며 요청에 응답합니다.
    \begin{examplebox}
    \textbf{예시:} 소셜 미디어 서비스. 일부 서버가 다운되어도 사용자는 계속해서 게시물을 보거나 올릴 수 있습니다.
    \end{examplebox}
    \item \textbf{Soft state (유연한 상태):} 시스템의 상태가 즉시 일관되지 않을 수 있습니다. 여러 노드에 데이터가 분산되어 저장될 때, 모든 노드가 동시에 최신 상태를 반영하지 못할 수 있습니다.
    \item \textbf{Eventual consistency (최종 일관성):} 데이터가 모든 노드에 즉시 전파되지 않더라도, 충분한 시간이 지나면 모든 노드의 데이터가 최종적으로 일관된 상태가 됩니다.
    \begin{examplebox}
    \textbf{예시:} 온라인 쇼핑몰의 '좋아요' 카운트. 사용자가 '좋아요'를 눌렀을 때, 즉시 모든 사용자에게 카운트가 반영되지 않을 수 있지만, 잠시 후에는 정확한 카운트가 표시됩니다. 몇 초의 지연은 허용됩니다.
    \end{examplebox}
\end{itemize}

\textbf{주요 사용 사례:}
\begin{itemize}
    \item 스키마가 자주 변경되거나 유연한 데이터 구조가 필요한 시나리오
    \item 대규모 웹 서비스, 소셜 미디어, IoT 데이터, 실시간 분석 등 고가용성과 확장성이 중요한 경우
\end{itemize}

\subsubsection*{8.3. CAP 이론 (Consistency, Availability, Partition tolerance)}

\begin{summarybox}
CAP 이론은 분산 컴퓨팅 시스템이 \textbf{일관성(Consistency)}, \textbf{가용성(Availability)}, \textbf{분할 내성(Partition tolerance)} 세 가지 속성 중 \textbf{동시에 두 가지만 만족할 수 있다}는 이론입니다. 시스템 설계 시 어떤 속성을 우선할지 선택해야 합니다.
\end{summarybox}

\textbf{CAP 속성:}
\begin{itemize}
    \item \textbf{Consistency (일관성):} 모든 클라이언트가 동일한 데이터를 보게 됩니다. 데이터가 업데이트되면 모든 노드에서 즉시 반영됩니다.
    \item \textbf{Availability (가용성):} 모든 요청에 대해 시스템이 항상 응답을 반환합니다. (응답이 성공적이지 않을 수 있어도)
    \item \textbf{Partition tolerance (분할 내성):} 네트워크 분할(일부 노드 간 통신 단절)이 발생하더라도 시스템이 계속 작동합니다.
\end{itemize}

\textbf{CAP 이론에 따른 시스템 분류:}
\begin{itemize}
    \item \textbf{CP 시스템 (Consistency + Partition tolerance):} 네트워크 분할 시 일관성을 유지하기 위해 가용성을 포기합니다. (예: 일부 관계형 데이터베이스, ZooKeeper)
    \item \textbf{AP 시스템 (Availability + Partition tolerance):} 네트워크 분할 시 가용성을 유지하기 위해 일관성을 포기합니다. (예: NoSQL 데이터베이스 - Cassandra, DynamoDB)
    \item \textbf{CA 시스템 (Consistency + Availability):} 분할 내성을 포기합니다. 단일 노드 시스템이나 네트워크 분할이 발생하지 않는 환경에서만 가능합니다. (예: 전통적인 단일 서버 관계형 데이터베이스)
\end{itemize}

\begin{warningbox}
실제로는 세 가지 속성을 모두 완벽하게 달성하기는 매우 어렵습니다. 특히 데이터 규모가 커질수록 데이터 일관성을 유지하는 것은 큰 도전 과제입니다.
\end{warningbox}

\subsection*{9. OLTP 시스템 vs OLAP 시스템}

데이터베이스 시스템은 그 목적에 따라 크게 두 가지로 나뉩니다.

\begin{adjustbox}{max width=\textwidth}
\begin{tabular}{lll}
\toprule
\textbf{구분} & \textbf{OLTP (Online Transactional Processing)} & \textbf{OLAP (Online Analytical Processing)} \\
\midrule
\textbf{목적} & 실시간 트랜잭션 처리 & 대규모 데이터 분석 및 보고 \\
\textbf{특징} &
\begin{itemize}
    \item 짧고 빈번한 트랜잭션 (읽기, 쓰기, 업데이트, 삭제)
    \item 높은 동시성 요구
    \item 데이터 일관성(ACID) 중요
    \item 정규화된 스키마
\end{itemize} &
\begin{itemize}
    \item 복잡하고 장기 실행되는 쿼리
    \item 대량의 데이터 읽기 중심
    \item 데이터 일관성보다 가용성/성능 중요 (BASE)
    \item 비정규화된 스키마 (스타 스키마, 스노우플레이크 스키마)
\end{itemize} \\
\textbf{기술 예시} &
\begin{itemize}
    \item \textbf{SQL 기반:} Oracle, SQL Server, MySQL
    \item \textbf{NoSQL 기반:} MongoDB (문서), Cassandra (컬럼), DynamoDB (키-값)
\end{itemize} &
\begin{itemize}
    \item \textbf{초기:} 데이터 웨어하우스, 데이터 마트
    \item \textbf{중기:} Hadoop
    \item \textbf{현대:} Databricks, Snowflake (클라우드 데이터 플랫폼)
\end{itemize} \\
\textbf{사용 사례} &
\begin{itemize}
    \item 온라인 상점의 주문 처리
    \item 은행의 계좌 거래
    \item 항공권 예약 시스템
\end{itemize} &
\begin{itemize}
    \item 월별 판매 보고서 생성
    \item 고객 행동 패턴 분석
    \item 시장 트렌드 예측
\end{itemize} \\
\bottomrule
\end{tabular}
\end{adjustbox}

\begin{figure}[h!]
    \centering
    % Placeholder for the diagram described in the transcript
    % The transcript describes a diagram with data producers -> ODS (OLTP) -> Data Warehouses/Data Marts/Data Lakes (OLAP) -> Serving Layer -> Data Consumers.
    % It also mentions Reverse ETL from OLAP back to OLTP.
    \begin{tcolorbox}[colback=gray!10, colframe=gray!50!black, title={OLTP 및 OLAP 시스템의 데이터 흐름 (개념도)}]
        \centering
        \textbf{데이터 생산자 (Data Producers)} \\
        (고객, 센서, 애플리케이션 등) \\
        $\downarrow$ \\
        \textbf{운영 데이터 저장소 (ODS: Operational Data Store)} \\
        (OLTP 시스템, 실시간 트랜잭션 처리) \\
        $\downarrow$ \\
        \textbf{데이터 웨어하우스 / 데이터 마트 / 데이터 레이크} \\
        (OLAP 시스템, 분석 및 보고) \\
        $\downarrow$ \\
        \textbf{서빙 계층 (Serving Layer)} \\
        (고성능 쿼리, 특정 애플리케이션용) \\
        $\downarrow$ \\
        \textbf{데이터 소비자 (Data Consumers)} \\
        (BI 분석가, 데이터 과학자, 애플리케이션 사용자) \\
        \textbf{역 ETL (Reverse ETL):} 분석 시스템에서 운영 시스템으로 통찰력 다시 주입
    \end{tcolorbox}
    \caption{OLTP와 OLAP 시스템 간의 데이터 흐름}
    \label{fig:oltp_olap_flow}
\end{figure}

\subsection*{10. 현대 데이터 플랫폼의 진화 과정}

현대 데이터 플랫폼은 다양한 데이터 소스, 유형, 처리 방식을 통합하여 비즈니스 사용 사례를 지원하도록 발전해 왔습니다.

\begin{adjustbox}{max width=\textwidth}
\begin{tabular}{lll}
\toprule
\textbf{영역} & \textbf{과거} & \textbf{현재} \\
\midrule
\textbf{데이터 소스} & 적은 수의 정형 데이터 소스 & 더 많은 데이터 소스, 고객에 대한 전체적인 시야 \\
\textbf{데이터 유형} & 주로 정형 데이터 (행/열, 스키마 명확) & 정형, 반정형 (JSON/XML), 비정형 (텍스트, 이미지, 비디오) \\
\textbf{데이터 수집} & 배치 처리 (일괄) & 배치, 실시간 스트리밍, 마이크로 배치 \\
\textbf{데이터 처리 (ETL)} & 데이터 웨어하우스 (구형) & 현대 데이터 플랫폼 (Snowflake, Databricks) \\
\textbf{사용 사례} & 보고서 (Reporting) & 보고서, 임시 쿼리 (EDA), 머신러닝 (ML) \\
\bottomrule
\end{tabular}
\end{adjustbox}

\subsection*{11. Hadoop과 Spark의 비교}

Hadoop은 데이터 레이크 시대의 시작을 알렸지만, Spark는 그 한계를 극복하고 현대 빅데이터 처리의 핵심 엔진이 되었습니다.

\subsubsection*{11.1. Hadoop의 특징}

\begin{summarybox}
Hadoop은 대규모 데이터를 저렴한 비용으로 분산 처리하기 위해 설계된 오픈소스 프레임워크입니다. HDFS(분산 파일 시스템)와 MapReduce(분산 처리 모델)가 핵심입니다.
\end{summarybox}

\textbf{장점:}
\begin{itemize}
    \item \textbf{저비용:} 값싼 범용 하드웨어(commodity grade hardware)를 사용하여 고가의 독점 서버를 대체할 수 있습니다.
    \item \textbf{신뢰성:} 데이터의 3중 복제(triple replication)를 통해 데이터 손실을 방지합니다.
    \item \textbf{수평적 확장성:} 필요에 따라 노드를 추가하여 용량과 성능을 쉽게 확장할 수 있습니다.
    \item \textbf{다양한 데이터 유형 지원:} 정형, 반정형, 비정형 데이터를 모두 저장하고 처리할 수 있습니다.
    \item \textbf{오픈소스 생태계:} Hive, Presto 등 다양한 오픈소스 프로젝트가 Hadoop을 기반으로 발전했습니다.
\end{itemize}

\textbf{한계:}
\begin{itemize}
    \item \textbf{복잡한 관리:} NameNode, YARN 등 여러 구성 요소를 관리하기 어렵습니다.
    \item \textbf{데이터 품질/신뢰성 문제:} 데이터 레이크의 유연성 때문에 데이터 거버넌스가 어려워질 수 있습니다.
    \item \textbf{느린 처리 속도:} MapReduce는 각 단계마다 디스크 I/O를 수행하여 느립니다.
    \item \textbf{잦은 업데이트에 부적합:} '한 번 쓰고 여러 번 읽기(write once, read many)'에 최적화되어 잦은 데이터 업데이트에는 비효율적입니다.
\end{itemize}

\subsubsection*{11.2. Spark의 특징}

\begin{summarybox}
Spark는 Hadoop의 MapReduce를 대체하는 빠르고 범용적인 분산 처리 엔진입니다. 인메모리(in-memory) 컴퓨팅과 DAG(Directed Acyclic Graph)를 활용하여 MapReduce보다 훨씬 빠른 속도를 제공합니다.
\end{summarybox}

\textbf{장점:}
\begin{itemize}
    \item \textbf{빠른 속도:} 데이터를 메모리에 캐싱하고 DAG를 사용하여 디스크 I/O를 최소화함으로써 Hadoop MapReduce보다 100배 이상 빠릅니다.
    \item \textbf{통합 분석 엔진:} 배치 처리, 스트리밍, SQL 쿼리, 머신러닝, 그래프 처리 등 다양한 워크로드를 하나의 엔진에서 지원합니다.
    \item \textbf{다국어 지원:} Scala(원래 언어), Python, Java, R 등 다양한 프로그래밍 언어를 지원합니다.
    \item \textbf{쉬운 배포:} Hadoop과 독립적으로 실행 가능하며, 다양한 클러스터 관리자(YARN, Mesos, Kubernetes)와 클라우드 환경에서 쉽게 배포됩니다.
    \item \textbf{수평적 확장성:} Hadoop과 마찬가지로 노드를 추가하여 쉽게 확장할 수 있습니다.
\end{itemize}

\textbf{Spark 아키텍처 (Hadoop과 유사):}
\begin{itemize}
    \item \textbf{Driver Program (드라이버 프로그램):} Spark 애플리케이션의 메인 함수를 실행하고, 클러스터의 Worker 노드에 작업을 분배합니다.
    \item \textbf{Cluster Manager (클러스터 관리자):} Worker 노드를 관리하고 리소스를 할당합니다.
    \item \textbf{Worker Nodes (워커 노드):} Driver로부터 작업을 받아 실행하며, Executor를 통해 CPU와 메모리를 사용합니다.
    \item \textbf{Executor (실행자):} Worker 노드에서 실제 작업을 수행하는 프로세스입니다.
\end{itemize}

\subsubsection*{11.3. Hadoop vs Spark 비교}

\begin{adjustbox}{max width=\textwidth}
\begin{tabular}{lll}
\toprule
\textbf{특성} & \textbf{Hadoop (MapReduce)} & \textbf{Spark} \\
\midrule
\textbf{처리 방식} & 디스크 기반 (각 Map/Reduce 단계마다 디스크 I/O) & 인메모리 기반 (DAG를 통해 디스크 I/O 최소화) \\
\textbf{속도} & 느림 & 빠름 (MapReduce 대비 100배 이상) \\
\textbf{데이터 흐름} & 선형 (Map $\rightarrow$ Shuffle $\rightarrow$ Reduce) & DAG (Directed Acyclic Graph) 기반의 유연한 처리 \\
\textbf{데이터 재사용} & 각 작업 후 디스크에 저장, 재사용 어려움 & 메모리에 RDD/DataFrame 캐싱, 재사용 용이 \\
\textbf{지원 언어} & Java (주로) & Scala, Python, Java, R 등 다국어 지원 \\
\textbf{사용 사례} & 대규모 배치 처리, 데이터 웨어하우징 & 배치, 스트리밍, SQL, ML, 그래프 처리 등 범용 \\
\textbf{복잡성} & 여러 컴포넌트 관리 (HDFS, YARN, MapReduce) & 단일 통합 엔진, 상대적으로 단순한 관리 \\
\bottomrule
\end{tabular}
\end{adjustbox}

\begin{notebox}
\textbf{RDD (Resilient Distributed Dataset):} Spark의 기본 데이터 추상화 단위입니다. 불변(immutable)하며, 분산된 데이터 컬렉션을 나타냅니다.
\textbf{DataFrame (데이터프레임):} RDD 위에 구축된 고수준 추상화로, 구조화된 데이터를 다루기 쉽고 최적화된 성능을 제공합니다. SQL과 유사한 방식으로 데이터를 조작할 수 있습니다.
\end{notebox}

\subsection*{12. 클라우드 컴퓨팅 모델 (IaaS, PaaS, SaaS)}

클라우드 서비스는 제공하는 추상화 수준에 따라 세 가지 모델로 나뉩니다.

\begin{adjustbox}{max width=\textwidth}
\begin{tabular}{llll}
\toprule
\textbf{모델} & \textbf{설명} & \textbf{관리 책임} & \textbf{예시} \\
\midrule
\textbf{IaaS (Infrastructure as a Service)} & 가상 머신, 스토리지, 네트워크 등 기본 인프라 제공 & 사용자가 OS, 미들웨어, 애플리케이션 관리 & AWS EC2, Azure Virtual Machines \\
\textbf{PaaS (Platform as a Service)} & 애플리케이션 개발 및 배포를 위한 플랫폼 제공 & 사용자가 애플리케이션만 관리, 인프라/OS는 클라우드 제공자 관리 & AWS Elastic Beanstalk, Azure App Service \\
\textbf{SaaS (Software as a Service)} & 완제품 소프트웨어를 서비스 형태로 제공 & 클라우드 제공자가 모든 것을 관리, 사용자는 소프트웨어 사용 & Gmail, Salesforce, Databricks \\
\bottomrule
\end{tabular}
\end{adjustbox}

\begin{examplebox}
\textbf{Databricks는 SaaS의 좋은 예시입니다.}
Databricks는 AWS, Azure, GCP와 같은 클라우드 인프라 위에 구축된 소프트웨어 서비스입니다. 사용자는 Databricks 플랫폼에 접속하여 데이터 파이프라인을 구축하고 분석을 수행할 수 있으며, 하부 인프라(서버, OS 등) 관리에 대해 걱정할 필요가 없습니다. 이는 사용자가 핵심 비즈니스 로직에 집중할 수 있게 해줍니다.
\end{examplebox}

\subsection*{13. 빅데이터 생태계의 주요 도전 과제}

빅데이터를 효과적으로 활용하는 데에는 여러 가지 어려움이 따릅니다.

\begin{itemize}
    \item \textbf{데이터 품질 및 신선도 (Data Quality & Staleness):}
    \begin{itemize}
        \item 데이터가 부정확하거나 불완전하면 잘못된 통찰로 이어집니다.
        \item 데이터가 너무 오래되면 그 가치가 떨어집니다 (데이터 온도 개념).
    \end{itemize}
    \item \textbf{데이터 사일로 (Data Silos):}
    \begin{itemize}
        \item 조직 내 여러 부서가 각자의 데이터를 독립적으로 관리하여 통합된 시야를 확보하기 어렵습니다.
        \item 서로 다른 기술 스택을 사용하여 데이터 공유 및 활용이 어렵습니다.
    \end{itemize}
    \item \textbf{파편화된 도구 (Fragmented Tools):}
    \begin{itemize}
        \item 다양한 클라우드 서비스, BI 도구, ML 프레임워크가 혼재하여 호환성 문제가 발생하고 생산성을 저해합니다.
    \end{itemize}
    \item \textbf{기술 부족 (Shortage of Skills):}
    \begin{itemize}
        \item 데이터를 다루고 최신 도구를 활용할 수 있는 숙련된 인력이 부족합니다.
    \end{itemize}
\end{itemize}
이러한 도전 과제들은 결국 부정확한 통찰과 지연된 결과로 이어져 데이터의 가치를 떨어뜨립니다.

\subsection*{14. 빅데이터 플랫폼을 위한 모범 사례}

효과적인 빅데이터 플랫폼을 구축하고 운영하기 위한 몇 가지 핵심 원칙입니다.

\begin{itemize}
    \item \textbf{결합도 낮은 시스템 구축 (Decoupled System):}
    \begin{itemize}
        \item 스토리지와 컴퓨팅 자원을 분리하여 각각 독립적으로 확장하고 관리할 수 있도록 합니다. (예: S3와 Spark)
    \end{itemize}
    \item \textbf{서비스 지향적 설계 (Service-Oriented):}
    \begin{itemize}
        \item 재사용 가능한 모듈 형태로 파이프라인이나 애플리케이션을 구축하여 특정 기능에 집중하도록 합니다.
    \end{itemize}
    \item \textbf{클라우드 스토리지 및 개방형 형식 활용 (Cloud Storage & Open Format):}
    \begin{itemize}
        \item 클라우드의 저렴하고 확장 가능한 스토리지를 활용하고, Parquet, Delta Lake와 같은 개방형 데이터 형식을 사용하여 특정 벤더에 종속되는 것을 피합니다 (Vendor Lock-in 방지).
    \end{itemize}
    \item \textbf{사용 사례에 맞는 도구 선택 (Right Tool for the Right Job):}
    \begin{itemize}
        \item 모든 문제에 하나의 도구가 최선일 수는 없습니다. 지연 시간, 처리량, 접근 패턴 등 사용 사례의 요구사항을 고려하여 최적의 도구를 선택합니다.
    \end{itemize}
    \item \textbf{로그 중심 설계 패턴 (Log-centric Design Patterns):}
    \begin{itemize}
        \item 변경 불가능한(immutable) 로그를 사용하여 데이터 변경 이력을 추적하고, 필요 시 데이터를 재처리(replay)할 수 있도록 합니다. (메달리온 아키텍처의 Bronze 계층의 중요성)
    \end{itemize}
    \item \textbf{소비자 요구에 따른 다양한 데이터 뷰 (Multiple Views of Data):}
    \begin{itemize}
        \item 동일한 원시 데이터로부터 다양한 소비자(보고서, ML 모델, 애플리케이션)의 요구에 맞춰 여러 형태의 뷰를 제공합니다. (예: 마스킹된 데이터, 집계된 데이터)
    \end{itemize}
    \item \textbf{구축 vs 구매 결정 (Build vs Buy):}
    \begin{itemize}
        \item 시간, 비용, 내부 전문성, 시장 경쟁 상황 등을 고려하여 솔루션을 직접 구축할지(Build) 외부 솔루션을 구매할지(Buy) 결정합니다.
    \end{itemize}
\end{itemize}

\subsection*{15. 기술 투자에 대한 비즈니스 정당성}

기술 투자는 단순히 비용이 아니라 비즈니스 성장을 위한 혁신 센터로 간주되어야 합니다.

\begin{itemize}
    \item \textbf{비즈니스 정당성 필수:} 새로운 기술을 도입할 때는 항상 명확한 비즈니스 목표와 정당성이 있어야 합니다. 단순히 '새로운 기술이라서' 도입하는 것은 실패할 가능성이 높습니다.
    \item \textbf{ROI (투자 수익률) 측정:} 모든 기술 투자는 측정 가능하고 가시적인 ROI를 가져와야 합니다. 핵심 성과 지표(KPI)를 설정하고 이를 통해 투자의 효과를 입증해야 합니다.
    \textbf{Capex vs Opex:}
    \begin{itemize}
        \item \textbf{Capex (Capital Expenditure):} 자본 지출. 온프레미스 서버 구매와 같이 초기 대규모 투자가 필요합니다.
        \item \textbf{Opex (Operational Expenditure):} 운영 지출. 클라우드 서비스 사용료와 같이 사용한 만큼 지불하는 방식입니다. 클라우드는 Opex 모델을 통해 초기 비용 부담을 줄여줍니다.
    \end{itemize}
    \item \textbf{데이터 거버넌스 및 보호:} 데이터는 중요한 자산이므로, 부적절한 접근과 유출로부터 보호되어야 합니다. Unity Catalog와 같은 도구를 사용하여 접근 권한을 세밀하게 제어해야 합니다.
    \item \textbf{기술적 관점과 비즈니스 관점의 통합:} 현재 상태(Current State)와 목표 상태(Desired State)를 명확히 하고, 기술적 해결책이 비즈니스 문제를 어떻게 해결할지 설명해야 합니다.
    \item \textbf{MEDDPICC 프레임워크 (Databricks 활용):}
    \begin{itemize}
        \item \textbf{M}etrics (측정 지표): KPI
        \item \textbf{E}conomic Buyer (경제적 구매자): 누가 비용을 지불할 것인가?
        \item \textbf{D}ecision Criteria (결정 기준): 지연 시간, 비용, 복원력 등
        \item \textbf{D}ecision Process (결정 과정): 누구와 협력해야 하는가?
        \item \textbf{P}ain (문제점): 현재 해결되지 않는 비즈니스 문제
        \item \textbf{I}mplicate (영향): 문제 해결 시 비즈니스에 미칠 긍정적 영향
        \item \textbf{C}hampion (옹호자): 프로젝트를 지지하는 내부 인물
        \item \textbf{C}ompetition (경쟁): 경쟁사 대비 우위 확보
    \end{itemize}
    이 프레임워크는 기술 투자의 비즈니스 가치를 명확히 하는 데 도움이 됩니다.
\end{itemize}

\newpage

% --- 절차/방법 (Procedures/Methods) ---
\section*{절차/방법: Databricks Free Edition 시작하기 (Lab 00)}
\addcontentsline{toc}{section}{절차/방법: Databricks Free Edition 시작하기 (Lab 00)}

Lab 00은 Databricks Free Edition 환경에서 기본적인 데이터 처리 작업을 수행하는 방법을 소개합니다.

\subsection*{1. Databricks Free Edition 계정 생성 및 접속}

\begin{enumerate}
    \item Databricks Free Edition 웹사이트에 접속하여 계정을 생성합니다. (일반 Trial 계정과 혼동하지 않도록 주의)
    \item 로그인 후, 화면 좌측 상단에 "Databricks Free Edition" 문구가 표시되는지 확인합니다. 이 문구가 없으면 Trial 계정일 가능성이 높습니다.
\end{enumerate}

\begin{warningbox}
\textbf{Trial 계정 vs Free Edition 계정:}
Databricks는 두 가지 유형의 무료 계정을 제공합니다.
\begin{itemize}
    \item \textbf{Trial Account (평가판 계정):} 14일 동안 모든 기능을 사용할 수 있지만, 이후에는 유료 전환이 필요합니다.
    \item \textbf{Free Edition Account (무료 에디션 계정):} 학술적 목적으로 영구적으로 무료이며, 일부 기능 제한이 있습니다. 이 수업에서는 Free Edition을 사용합니다.
\end{itemize}
만약 "Databricks Free Edition" 문구가 보이지 않는다면, Trial 계정일 수 있으므로 다시 Free Edition 계정을 생성해야 합니다.
\end{warningbox}

\subsection*{2. Lab 00 노트북 가져오기 (Import)}

\begin{enumerate}
    \item 제공된 Google Drive 링크에서 Lab 00 zip 파일을 다운로드합니다.
    \item Databricks Workspace 좌측 패널에서 \textbf{Workspace}를 클릭합니다.
    \item 본인의 사용자 이름 아래에 폴더를 생성합니다 (예: \texttt{labs}).
        \begin{itemize}
            \item 사용자 이름 클릭 $\rightarrow$ \textbf{Create} $\rightarrow$ \textbf{Folder} $\rightarrow$ 이름 입력 (예: \texttt{labs})
        \end{itemize}
    \item 생성된 폴더 안에서 마우스 오른쪽 버튼 클릭 $\rightarrow$ \textbf{Import}를 선택합니다.
    \item 다운로드한 zip 파일을 드래그 앤 드롭하여 가져옵니다. Databricks가 자동으로 압축을 해제하고 노트북 파일들을 생성합니다.
    \item 가져오기가 완료되면 \texttt{lab00} 폴더 안에 \texttt{readme}, \texttt{initialize}, \texttt{intro\_basics}, \texttt{read\_write}, \texttt{spark\_intro} 5개의 노트북이 보일 것입니다.
\end{enumerate}

\subsection*{3. 컴퓨트(Compute) 연결 및 노트북 실행}

Databricks Free Edition은 서버리스(Serverless) 환경이므로, 별도의 클러스터를 생성할 필요 없이 컴퓨트 자원에 연결하기만 하면 됩니다.

\begin{enumerate}
    \item \texttt{initialize} 노트북을 엽니다.
    \item 노트북 상단의 \textbf{Connect} 버튼을 클릭하고 \textbf{Serverless}를 선택합니다.
    \item 연결이 완료되면, 첫 번째 셀을 실행합니다. (셀 좌측의 실행 버튼 클릭 또는 \texttt{Ctrl+Enter})
    \item 첫 번째 셀은 카탈로그, 스키마, 볼륨을 생성하는 명령을 포함합니다.
        \begin{lstlisting}[language=sql, caption={initialize 노트북의 첫 번째 셀 예시}, label={lst:initialize_cell1}]
CREATE CATALOG IF NOT EXISTS csci103_catalog;
USE CATALOG csci103_catalog;

CREATE SCHEMA IF NOT EXISTS lab00;
USE SCHEMA lab00;

CREATE VOLUME IF NOT EXISTS lab00_input;
CREATE VOLUME IF NOT EXISTS lab00_output;
        \end{lstlisting}
    \item 실행이 성공하면, 좌측 패널의 \textbf{Catalog} 아이콘을 클릭하여 생성된 \texttt{csci103\_catalog}와 그 아래의 \texttt{lab00} 스키마, 그리고 \texttt{lab00\_input}, \texttt{lab00\_output} 볼륨을 확인할 수 있습니다.
\end{enumerate}

\begin{notebox}
\textbf{Databricks의 메타데이터 계층 구조:}
\begin{itemize}
    \item \textbf{Catalog (카탈로그):} 최상위 메타데이터 컨테이너. 데이터베이스 서버와 유사합니다.
    \item \textbf{Schema (스키마):} 카탈로그 아래의 논리적 그룹. 전통적인 데이터베이스의 스키마 또는 데이터베이스와 유사합니다.
    \item \textbf{Table (테이블):} 스키마 아래에 저장되는 실제 데이터 테이블.
    \item \textbf{Volume (볼륨):} 클라우드 스토리지 내의 특정 경로를 가리키는 논리적 단위. 파일 시스템 경로를 추상화합니다.
\end{itemize}
\end{notebox}

\subsection*{4. Databricks 노트북 기본 사용법}

\begin{itemize}
    \item \textbf{셀 실행:} 셀 좌측의 실행 버튼 또는 \texttt{Ctrl+Enter}.
    \item \textbf{언어 전환 (Magic Commands):}
        \begin{itemize}
            \item \texttt{\%python}: Python 코드 실행 (기본값)
            \item \texttt{\%sql}: SQL 쿼리 실행
            \item \texttt{\%sh}: 셸(Shell) 명령어 실행 (예: \texttt{ls}, \texttt{top})
            \item \texttt{\%fs}: Databricks 파일 시스템(DBFS) 명령어 실행 (예: \texttt{ls}, \texttt{cp})
            \item \texttt{\%md}: Markdown 텍스트 작성
        \\end{itemize}
    \item \textbf{데이터셋 탐색:}
        \begin{itemize}
            \item \texttt{\%fs ls /databricks-datasets}: Databricks에서 제공하는 공개 데이터셋 목록 확인.
            \item \texttt{dbutils.fs.ls("경로")}: Python에서 파일 시스템 목록 확인.
            \item \texttt{display(dbutils.fs.ls("경로"))}: 결과를 표 형태로 예쁘게 출력.
        \end{itemize}
\end{itemize}

\begin{examplebox}
\textbf{다양한 언어 셀 예시:}
\begin{lstlisting}[language=python, caption={Python 셀 예시}]
print("Hello from Python!")
my_variable = "Python variable"
print(my_variable)
\end{lstlisting}

\begin{lstlisting}[language=sql, caption={SQL 셀 예시}]
%sql
SELECT "Hello from SQL!" AS greeting;
SELECT current_user() AS current_user;
\end{lstlisting}

\begin{lstlisting}[language=sh, caption={Shell 셀 예시}]
%sh
ls -l /databricks-datasets/
\end{lstlisting}

\begin{lstlisting}[language=markdown, caption={Markdown 셀 예시}]
%md
# Markdown 제목
**굵은 글씨**와 *기울임 글씨*를 사용할 수 있습니다.
- 목록 1
- 목록 2
\end{lstlisting}
\end{examplebox}

\subsection*{5. 데이터프레임 생성 및 조작}

Spark의 핵심은 데이터프레임을 사용하여 데이터를 처리하는 것입니다.

\begin{enumerate}
    \item \textbf{데이터 읽기:}
        \begin{lstlisting}[language=python, caption={CSV 파일 읽기 예시}, label={lst:read_csv}]
file_path = "/databricks-datasets/bikeSharing/data-001/201508_trip_data.csv"
df = spark.read.format("csv") \
  .option("header", "true") \
  .option("inferSchema", "true") \
  .load(file_path)
df.printSchema()
df.show(5)
        \end{lstlisting}
        \begin{itemize}
            \item \texttt{format("csv")}: 파일 형식 지정.
            \item \texttt{option("header", "true")}: 첫 줄을 헤더로 인식.
            \item \texttt{option("inferSchema", "true")}: Spark가 데이터 타입을 자동으로 추론. (정확한 스키마를 알고 있다면 수동으로 정의하는 것이 더 좋습니다.)
            \item \texttt{printSchema()}: 데이터프레임의 스키마(컬럼 이름과 타입) 출력.
            \item \texttt{show(n)}: 데이터프레임의 상위 n개 행 출력.
        \end{itemize}
    \item \textbf{스키마 수동 정의:}
        \begin{lstlisting}[language=python, caption={스키마 수동 정의 예시}]
from pyspark.sql.types import StructType, StructField, StringType, IntegerType

custom_schema = StructType([
    StructField("trip_id", IntegerType(), True),
    StructField("duration", IntegerType(), True),
    StructField("start_date", StringType(), True),
    # ... 다른 컬럼들
])
df_custom = spark.read.format("csv") \
  .option("header", "true") \
  .schema(custom_schema) \
  .load(file_path)
df_custom.printSchema()
        \end{lstlisting}
        \begin{itemize}
            \item \texttt{schema(custom_schema)}: Spark의 스키마 추론 대신 사용자가 정의한 스키마를 적용.
        \end{itemize}
    \item \textbf{데이터 쓰기 (다양한 형식):}
        \begin{lstlisting}[language=python, caption={데이터를 Parquet 형식으로 쓰기 예시}]
output_path_parquet = "/Volumes/csci103_catalog/lab00/lab00_output/bike_data_parquet"
df.write.format("parquet").mode("overwrite").save(output_path_parquet)
        \end{lstlisting}
        \begin{itemize}
            \item \texttt{mode("overwrite")}: 이미 파일이 존재하면 덮어씁니다. (\texttt{"append"}, \texttt{"ignore"}, \texttt{"errorifexists"} 등)
            \item \texttt{save(경로)}: 지정된 경로에 데이터를 저장합니다.
            \item \textbf{Delta Lake 형식으로 쓰기:} Databricks의 기본이자 권장 형식입니다.
                \begin{lstlisting}[language=python, caption={데이터를 Delta 형식으로 쓰기 예시}]
output_path_delta = "/Volumes/csci103_catalog/lab00/lab00_output/bike_data_delta"
df.write.format("delta").mode("overwrite").save(output_path_delta)
                \end{lstlisting}
        \end{itemize}
    \item \textbf{테이블 생성 및 등록 (Unity Catalog):}
        \begin{lstlisting}[language=sql, caption={Delta 테이블 생성 및 등록 예시}]
%sql
CREATE TABLE IF NOT EXISTS bike_trips_delta
USING DELTA
LOCATION '/Volumes/csci103_catalog/lab00/lab00_output/bike_data_delta';

-- 또는 데이터프레임을 사용하여 테이블로 저장
-- df.write.format("delta").mode("overwrite").saveAsTable("bike_trips_df_table")
        \end{lstlisting}
        \begin{itemize}
            \item \texttt{CREATE TABLE ... USING DELTA LOCATION ...}: 기존 Delta 형식으로 저장된 데이터를 Unity Catalog에 테이블로 등록합니다.
            \item \texttt{saveAsTable("테이블_이름")}: 데이터프레임을 직접 Unity Catalog에 테이블로 등록합니다. 이 경우 Databricks가 스토리지 경로를 관리합니다 (Managed Table).
            \item 테이블이 등록되면 \textbf{Catalog Explorer}에서 해당 테이블의 스키마, 히스토리, 권한 등을 확인할 수 있습니다.
        \end{itemize}
    \item \textbf{컬럼 추가/삭제/변환:}
        \begin{lstlisting}[language=python, caption={데이터프레임 컬럼 조작 예시}]
from pyspark.sql.functions import col, lit

# 새로운 컬럼 추가
df_with_new_col = df.withColumn("new_column", lit("some_value"))

# 기존 컬럼 이름 변경
df_renamed = df_with_new_col.withColumnRenamed("duration", "trip_duration_seconds")

# 컬럼 삭제
df_dropped = df_renamed.drop("start_station_id", "end_station_id")

df_dropped.printSchema()
        \end{lstlisting}
    \item \textbf{사용자 정의 함수 (UDF) 사용:}
        \begin{lstlisting}[language=python, caption={사용자 정의 함수 (UDF) 예시}]
from pyspark.sql.functions import udf
from pyspark.sql.types import IntegerType

# duration을 분 단위로 변환하는 UDF 정의
def seconds_to_minutes(seconds):
    return seconds // 60

seconds_to_minutes_udf = udf(seconds_to_minutes, IntegerType())

# UDF를 사용하여 새로운 컬럼 생성
df_with_minutes = df.withColumn("trip_duration_minutes", seconds_to_minutes_udf(col("duration")))
df_with_minutes.select("duration", "trip_duration_minutes").show(5)
        \end{lstlisting}
\end{enumerate}

\newpage

% --- 실습/코드/오류와 해결 (Lab/Code/Troubleshooting) ---
\section*{실습/코드/오류와 해결}
\addcontentsline{toc}{section}{실습/코드/오류와 해결}

\subsection*{1. Databricks UI 탐색 팁}

\begin{itemize}
    \item \textbf{좌측 패널:} Workspace (노트북, 폴더), Catalog (데이터, 테이블, 볼륨), Compute (클러스터 관리), Workflow (작업 스케줄링) 등 주요 기능에 접근합니다.
    \item \textbf{노트북 상단:} 언어 선택 (Python, SQL 등), 컴퓨트 연결/해제, 셀 실행 옵션 (모든 셀 실행, 선택 셀 실행 등).
    \item \textbf{셀 출력:} 셀 실행 후 결과는 셀 아래에 표시됩니다. 표, 그래프, 텍스트 등 다양한 형태로 나타납니다.
    \item \textbf{쿼리 프로파일 (Query Profile):} Spark 쿼리 실행 후 셀 하단의 'Performance' 탭을 클릭하면 쿼리 실행 계획, 병목 현상 등을 시각적으로 확인할 수 있습니다. 이는 Spark UI의 Databricks 버전입니다.
\end{itemize}

\subsection*{2. Lab 00에서 발생할 수 있는 일반적인 오류 및 해결책}

\begin{adjustbox}{max width=\textwidth}
\begin{tabular}{lll}
\toprule
\textbf{증상} & \textbf{원인 가설} & \textbf{수정 방법} \\
\midrule
\texttt{Notebook failed to load} &
\begin{itemize}
    \item Free Edition이 아닌 Trial 계정 사용
    \item 네트워크 문제 또는 Databricks 서비스 일시적 오류
\end{itemize} &
\begin{itemize}
    \item Databricks Free Edition 계정으로 다시 생성 및 로그인
    \item 인터넷 연결 확인 후 새로고침 또는 잠시 후 재시도
\end{itemize} \\
\texttt{Permission denied} 또는 \texttt{Catalog not found} &
\begin{itemize}
    \item \texttt{initialize} 노트북을 실행하지 않음
    \item 다른 사용자의 카탈로그/스키마에 접근 시도
    \item Free Edition이 아닌 Trial 계정 사용
\end{itemize} &
\begin{itemize}
    \item \texttt{initialize} 노트북을 먼저 실행하여 카탈로그/스키마/볼륨 생성
    \item 본인의 카탈로그/스키마를 사용하거나, 다른 사용자에게 권한 요청
    \item Free Edition 계정으로 다시 생성 및 로그인
\end{itemize} \\
\texttt{Command failed} 또는 \texttt{Error in SQL statement} &
\begin{itemize}
    \item 구문 오류 (SQL, Python 등)
    \item 컴퓨트(Serverless)에 연결되지 않음
    \item 데이터 경로 오류
\end{itemize} &
\begin{itemize}
    \item 코드의 오타, 문법 오류 확인
    \item 노트북 상단의 \textbf{Connect} 버튼을 클릭하여 Serverless 컴퓨트에 연결
    \item 파일 경로가 올바른지 확인 (특히 \texttt{/Volumes/catalog/schema/volume} 형식)
\end{itemize} \\
\texttt{DBFS writing data to DBFS is not supported} &
\begin{itemize}
    \item Free Edition에서 DBFS 유틸리티를 통한 직접 쓰기 제한
\end{itemize} &
\begin{itemize}
    \item \texttt{/Volumes/csci103_catalog/lab00/lab00_output/bike_data_delta}와 같은 \texttt{/Volumes} 경로를 통해 데이터를 저장합니다.
    \item \texttt{dbutils.fs.put()} 대신 \texttt{df.write.save()} 또는 \texttt{df.write.saveAsTable()} 사용.
\end{itemize} \\
\bottomrule
\end{tabular}
\end{adjustbox}

\subsection*{3. Spark의 실행 계획 확인 (\texttt{explain()})}

Spark는 코드를 즉시 실행하지 않고, 최적화된 실행 계획을 세운 후 실행합니다. \texttt{explain()} 함수를 사용하여 이 계획을 확인할 수 있습니다.

\begin{lstlisting}[language=python, caption={Spark 실행 계획 확인 예시}, label={lst:spark_explain}]
# 이전 예시에서 생성된 데이터프레임 df_with_minutes 사용
df_with_minutes.explain(extended=True)
\end{lstlisting}
\begin{itemize}
    \item \texttt{explain()}: Spark가 데이터프레임 작업을 어떻게 수행할지 논리적/물리적 실행 계획을 보여줍니다.
    \item \texttt{extended=True}: 더 상세한 정보를 포함합니다.
    \item \textbf{Photon:} Databricks의 고성능 쿼리 엔진으로, Spark 작업을 가속화합니다. \texttt{explain()} 결과에서 Photon이 지원하는 작업과 지원하지 않는 작업을 확인할 수 있습니다.
\end{itemize}

\newpage

% --- 체크리스트 (Checklist) ---
\section*{체크리스트}
\addcontentsline{toc}{section}{체크리스트}

\subsection*{1. 과제 제출 전 점검표}

\begin{itemize}
    \item \textbf{마감 기한 준수:} 모든 과제는 마감일 자정까지 제출해야 합니다. (업로드 문제 대비 최소 5분 전 제출 권장)
    \item \textbf{지연 정책 확인:} 지연 제출 시 하루당 2점 감점, 최대 10점 감점됩니다. (늦더라도 제출하는 것이 중요)
    \item \textbf{개별 제출:} 모든 과제는 개별 제출입니다. (그룹 프로젝트는 별도 명시)
    \item \textbf{노트북 내보내기:} 완료된 노트북을 Canvas에 업로드할 수 있도록 내보내기(export)했는지 확인합니다.
    \item \textbf{초기 시작:} 과제는 가능한 한 빨리 시작하여 충분한 학습 시간을 확보합니다.
    \item \textbf{질문 활용:} 막히는 부분이 있다면 Slack, Office Hours를 통해 질문합니다. 마감 직전까지 기다리지 않습니다.
\end{itemize}

\subsection*{2. 학습 점검표}

\begin{itemize}
    \item 데이터 엔지니어링의 정의와 중요성을 설명할 수 있는가?
    \item 빅데이터의 5V 특성을 각각 설명하고 예시를 들 수 있는가?
    \item 데이터 페르소나(엔지니어, 과학자, 분석가 등)의 역할을 구분할 수 있는가?
    \item SQL(ACID)과 NoSQL(BASE) 시스템의 차이점과 사용 사례를 비교할 수 있는가?
    \item CAP 이론을 설명하고, 시스템 설계 시의 트레이드오프를 이해하는가?
    \item OLTP와 OLAP 시스템의 목적과 특징을 구분할 수 있는가?
    \item Hadoop과 Spark의 주요 차이점(속도, 처리 방식 등)을 설명할 수 있는가?
    \item IaaS, PaaS, SaaS 클라우드 모델의 차이점을 설명하고 예시를 들 수 있는가?
    \item Databricks Free Edition 환경에서 노트북 가져오기, 컴퓨트 연결, 셀 실행을 할 수 있는가?
    \item Databricks에서 카탈로그, 스키마, 볼륨을 생성하고 탐색할 수 있는가?
    \item Spark 데이터프레임을 사용하여 CSV, JSON, Parquet, Delta 파일을 읽고 쓸 수 있는가?
    \item Spark 데이터프레임에서 컬럼 추가/삭제/변환 및 UDF를 사용할 수 있는가?
\end{itemize}

\newpage

% --- FAQ (Frequently Asked Questions) ---
\section*{FAQ (자주 묻는 질문)}
\addcontentsline{toc}{section}{FAQ (자주 묻는 질문)}

\begin{itemize}
    \item \textbf{Q: 최종 발표는 개인 발표인가요, 그룹 발표인가요?}
    \textbf{A:} 최종 발표는 그룹 프로젝트의 일환으로 진행되는 그룹 발표입니다. 각 팀원은 데이터 엔지니어링, 데이터 과학 등 특정 페르소나 역할을 맡아 프로젝트의 해당 측면을 발표하게 됩니다. 이는 실제 업무 환경을 반영한 것입니다.

    \item \textbf{Q: 강의 섹션(Section)에도 반드시 참석해야 하나요?}
    \textbf{A:} 섹션은 과제 지원 및 강의 자료 복습에 중점을 둡니다. 라이브 참석이 권장되지만, 시간대 문제 등으로 참석이 어렵다면 녹화본이 제공되므로 반드시 시청해야 합니다. 특히 과제 리뷰가 있는 주에는 시청이 필수적입니다.

    \item \textbf{Q: 교재는 구매해야 하나요?}
    \textbf{A:} 두 권의 교재(\textit{Simplifying Data Engineering and Analytics with Delta}, \textit{Spark, the Definitive Guide}) 모두 하버드 도서관을 통해 전자책 형태로 무료로 이용할 수 있습니다. 온라인 구매도 가능하지만 필수는 아닙니다.

    \item \textbf{Q: Databricks Free Edition 계정 생성 시 오류가 발생했습니다.}
    \textbf{A:} "Notebook failed to load"와 같은 오류는 주로 Databricks Free Edition이 아닌 일반 Trial 계정을 생성했을 때 발생합니다. Free Edition 계정 생성 링크를 다시 확인하고, 계정 생성 후 Databricks UI 좌측 상단에 "Databricks Free Edition" 문구가 표시되는지 꼭 확인하세요.
\end{itemize}

\newpage

% --- 빠르게 훑어보기 (1페이지 요약) ---
\section*{빠르게 훑어보기 (1페이지 요약)}
\addcontentsline{toc}{section}{빠르게 훑어보기 (1페이지 요약)}

\begin{tcolorbox}[colback=myblue!20, colframe=blue!50!black, title=\textbf{CSCI103 Lecture 1 핵심 요약}, sharp corners, boxsep=5mm, left=3mm, right=3mm, top=3mm, bottom=3mm]

\begin{itemize}
    \item[] \textbf{데이터 엔지니어링:} 원시 데이터를 가치 있는 통찰로 변환하는 핵심 역할. (수집 $\rightarrow$ 변환 $\rightarrow$ 저장 $\rightarrow$ 제공)
    \item[] \textbf{빅데이터 5V:} \textbf{Volume} (규모), \textbf{Velocity} (속도), \textbf{Variety} (다양성), \textbf{Veracity} (정확성), \textbf{Value} (가치).
    \item[] \textbf{데이터 페르소나:} 엔지니어, 과학자, 분석가, DevOps, 리더 등 다양한 역할의 협업.
    \item[] \textbf{데이터 온도:} Hot (실시간), Warm (인터랙티브), Cold (배치) 데이터에 따라 처리 속도 결정.
    \item[] \textbf{데이터 플랫폼 진화:} DBMS $\rightarrow$ DW $\rightarrow$ Hadoop $\rightarrow$ Spark $\rightarrow$ \textbf{레이크하우스}.
    \item[] \textbf{SQL (ACID) vs NoSQL (BASE):}
    \begin{itemize}
        \item \textbf{SQL:} 일관성, 신뢰성 중시 (금융, 거래). ACID (원자성, 일관성, 격리성, 영속성).
        \item \textbf{NoSQL:} 유연성, 확장성, 가용성 중시 (웹, 소셜). BASE (기본적 가용성, 유연한 상태, 최종 일관성).
    \end{itemize}
    \item[] \textbf{CAP 이론:} 분산 시스템은 일관성(C), 가용성(A), 분할 내성(P) 중 \textbf{두 가지만} 만족 가능.
    \item[] \textbf{OLTP vs OLAP:}
    \begin{itemize}
        \item \textbf{OLTP:} 실시간 트랜잭션 처리 (예: 은행 거래).
        \item \textbf{OLAP:} 대규모 데이터 분석 및 보고 (예: 판매 보고서).
    \end{itemize}
    \item[] \textbf{Hadoop vs Spark:}
    \begin{itemize}
        \item \textbf{Hadoop:} 디스크 기반, 느림, 데이터 레이크 시작.
        \item \textbf{Spark:} 인메모리 기반, 빠름, 범용 분석 엔진 (배치, 스트리밍, ML, SQL).
    \end{itemize}
    \item[] \textbf{클라우드 모델:} \textbf{IaaS} (인프라), \textbf{PaaS} (플랫폼), \textbf{SaaS} (소프트웨어). Databricks는 SaaS.
    \item[] \textbf{Databricks Lab 00:}
    \begin{itemize}
        \item Free Edition 계정 생성 및 접속 확인.
        \item 노트북 가져오기, Serverless 컴퓨트 연결.
        \item 카탈로그, 스키마, 볼륨 생성 (\texttt{initialize} 노트북).
        \item \texttt{\%sql}, \texttt{\%python}, \texttt{\%sh}, \texttt{\%fs}, \texttt{\%md} 등 매직 명령어 활용.
        \item Spark 데이터프레임으로 데이터 읽기/쓰기 (CSV, JSON, Parquet, Delta).
        \item \texttt{CREATE TABLE} 또는 \texttt{saveAsTable()}로 테이블 등록.
        \item 컬럼 조작, UDF 사용, \texttt{explain()}으로 실행 계획 확인.
    \end{itemize}
    \item[] \textbf{모범 사례:} 스토리지/컴퓨트 분리, 오픈 형식, 사용 사례에 맞는 도구, 로그 중심 설계.
    \item[] \textbf{비즈니스 정당성:} 기술 투자는 ROI를 통해 비즈니스 가치를 입증해야 함 (MEDDPICC).
\end{itemize}
\end{tcolorbox}

\newpage

% --- 부록 (Appendix) ---
\section*{부록: 강의 운영 및 평가}
\addcontentsline{toc}{section}{부록: 강의 운영 및 평가}

\subsection*{1. 강의 운영 방식}

\begin{itemize}
    \item \textbf{강의 (Lectures):} 주 1회 진행되며, 이론 학습 (전반부)과 실습 (후반부)으로 구성됩니다. 라이브 Zoom 세션 또는 녹화본 시청이 가능합니다.
    \item \textbf{섹션 (Sections):} 목요일 오후 6시~7시 (미국 동부 시간)에 진행됩니다. 과제 지원, 강의 자료 복습, 과제 풀이 등에 중점을 둡니다. 녹화본이 제공됩니다.
    \item \textbf{오피스 아워 (Office Hours):} 교수진과의 1:1 세션으로, 과제나 학습에 어려움이 있을 때 도움을 받을 수 있습니다. Cali 링크를 통해 쉽게 예약 가능합니다.
    \item \textbf{주요 소통 채널:}
    \begin{itemize}
        \item \textbf{Canvas:} 강의 자료 (과제, 슬라이드, 강의 계획서, 일정), 녹화본 링크가 게시됩니다.
        \item \textbf{Slack:} 질문, 토론, 그룹 프로젝트 협업 등 상호작용적인 소통을 위한 채널입니다.
    \end{itemize}
\end{itemize}

\subsection*{2. 학생 기대 사항}

\begin{itemize}
    \item \textbf{적극적인 참여 및 질문:} 궁금한 점은 주저하지 말고 질문합니다.
    \item \textbf{개방적인 학습 태도:} 새로운 아이디어와 기술에 대해 배우려는 의지를 가집니다.
    \item \textbf{예의 바르고 전문적인 태도:} 교수진 및 다른 학생들에게 항상 예의를 지키고 존중합니다.
    \item \textbf{그룹 프로젝트 참여:} 팀원들과 적극적으로 협력하고, 자신의 역할을 다하여 최대한의 학습 효과를 얻습니다. 팀원들과 일찍 소통하여 관계를 구축하는 것이 중요합니다.
\end{itemize}

\subsection*{3. 평가 및 채점 기준}

\begin{adjustbox}{max width=\textwidth}
\begin{tabular}{ll}
\toprule
\textbf{평가 항목} & \textbf{설명} \\
\midrule
\textbf{과제 (Assignments)} & 총 4개의 과제가 주어집니다. 강의 내용과 연관된 실습 위주입니다. \\
\textbf{사례 연구 (Case Studies)} & 총 2개의 사례 연구가 있습니다. \\
\textbf{퀴즈 (Quizzes)} & 총 2개의 퀴즈가 있습니다. \\
\textbf{최종 발표 (Final Presentation)} & 그룹 프로젝트의 결과물을 발표합니다. \\
\textbf{참여도 (Participation)} & 강의 중 질문, Slack 활동, 그룹 프로젝트 참여도 등이 반영됩니다. \\
\bottomrule
\end{tabular}
\end{adjustbox}

\begin{warningbox}
\textbf{과제 마감 정책:}
\begin{itemize}
    \item 모든 과제는 마감일 자정(midnight)까지 제출해야 합니다.
    \item 업로드 문제에 대비하여 마감 시간 최소 5분 전에는 제출하는 것이 좋습니다.
    \item 지연 제출 시 하루당 2점 감점되며, 과제당 최대 10점까지 감점됩니다. (늦더라도 제출하면 부분 점수를 받을 수 있습니다.)
    \item 가능한 한 빨리 과제를 시작하여 충분한 시간을 가지고 학습하고 제출하는 것이 가장 좋은 방법입니다.
\end{itemize}
\end{warningbox}

\subsection*{4. 학습 로드맵 (초심자 권장)}

\begin{enumerate}
    \item \textbf{기초 다지기:} 데이터 엔지니어링의 정의, 빅데이터 5V, 데이터 페르소나 등 핵심 개념을 먼저 이해합니다.
    \item \textbf{환경 설정:} Databricks Free Edition 계정을 성공적으로 생성하고 Lab 00 노트북을 실행하여 기본적인 환경에 익숙해집니다.
    \item \textbf{데이터 처리 기본:} Spark 데이터프레임을 사용하여 데이터 읽기, 쓰기, 간단한 변환 작업을 반복적으로 실습합니다.
    \item \textbf{개념 심화:} 데이터 플랫폼의 진화, SQL/NoSQL, OLTP/OLAP, CAP 이론 등 이론적 배경을 학습하여 '왜' 이러한 기술이 필요한지 이해합니다.
    \item \textbf{과제 해결:} 각 과제를 통해 배운 이론과 실습 지식을 적용하고, 막히는 부분은 적극적으로 질문하여 해결합니다.
    \item \textbf{프로젝트 준비:} 그룹 프로젝트를 통해 실제 비즈니스 문제를 해결하는 과정을 경험하고, 다양한 데이터 페르소나의 역할을 수행해봅니다.
\end{enumerate}

\end{document}